% GENERATED FILE. Edit canonical lessons, then run npm run generate:books.
% canonical-source-hash: fnv1a64:34103a90654cd40d
% canonical-lessons: GU-C264-padoshi, GU-C264-dost, GU-C264-bahenpani, GU-C264-olakhvu, GU-C264-bhetvu

\chapter{Friends and neighbours}
\label{ch:gu-friends-and-neighbours}
\hlchaptermodality{264}

\begin{chapteropening}
\textbf{By the end of this chapter:} \emph{I can ask whether people know each other and say who is not at home.}

Complete the last lesson of chapter 264: I can ask whether people know each other and say who is not at home.
\end{chapteropening}

\section[\texorpdfstring{\gu{પડોશી} (paḍośī)}{પડોશી (paḍośī)}]{\gu{પડોશી} (paḍośī) — neighbour}

\label{lesson:GU-C264-padoshi}



\begin{quote}
Before the new one: say the Gujarati for mother-in-law, then the Gujarati for nephew.
\end{quote}

\subsection*{You'll want to know: \gu{પડોશી}}
\textbf{\gu{પડોશી}} — \emph{paḍośī} — \textquotedblleft{}neighbour\textquotedblright{}.

\begin{quote}
\textbf{\gu{અમારા} \gu{પડોશી} \gu{ઘરે} \gu{નથી}.} — \emph{amārā paḍośī ghare nathī} — \textquotedblleft{}Our neighbours are not at home.\textquotedblright{}
\end{quote}

\begin{cousinweb}
\textbf{\gu{પડોશી}} (\emph{paḍośī}) is a tadbhava. It comes from Sanskrit \emph{prativēśya}, \textquotedblleft{}neighbour\textquotedblright{}; CDIAL lists Gujarati \emph{paṛośī} (CDIAL 8599).
\end{cousinweb}

\begin{grammarlens}[title={What you've built}]
One more word for everyday life at A2.
\end{grammarlens}

\subsection*{Your turn}
\begin{itemize}
\raggedright
  \item \emph{Say it:} \emph{paḍośī}
  \item \emph{Say it:} \emph{paḍośī}, once more
  \item \emph{Say it:} \emph{amārā paḍośī ghare nathī} — Our neighbours are not at home.
  \item \emph{Recall:} say the Gujarati for mother-in-law, then the Gujarati for nephew, then say \emph{paḍośī} again
\end{itemize}

\begin{tcolorbox}[breakable,colback=teal!4,colframe=teal!35!black,title={Before you move on}]
What does \gu{પડોશી} mean? (\textquotedblleft{}Neighbour\textquotedblright{}.) Say it once more.
\end{tcolorbox}
\section[\texorpdfstring{\gu{દોસ્ત} (dost)}{દોસ્ત (dost)}]{\gu{દોસ્ત} (dost) — pal, close friend}

\label{lesson:GU-C264-dost}



\begin{quote}
Before the new one: say the Gujarati for nephew, then the Gujarati for neighbour.
\end{quote}

\subsection*{You'll want to know: \gu{દોસ્ત}}
\textbf{\gu{દોસ્ત}} — \emph{dost} — \textquotedblleft{}pal, close friend\textquotedblright{}.

\begin{quote}
\textbf{\gu{શું} \gu{તે} \gu{તમારો} \gu{દોસ્ત} \gu{છે}?} — \emph{śũ te tamāro dost che?} — \textquotedblleft{}Is he your pal?\textquotedblright{}
\end{quote}

\begin{cousinweb}
\textbf{\gu{દોસ્ત}} (\emph{dost}) is a loan, usually traced to Persian \emph{dost}, \textquotedblleft{}friend\textquotedblright{}. CDIAL has no entry for it.
\end{cousinweb}

\begin{grammarlens}[title={What you've built}]
One more word for everyday life at A2.
\end{grammarlens}

\subsection*{Your turn}
\begin{itemize}
\raggedright
  \item \emph{Say it:} \emph{dost}
  \item \emph{Say it:} \emph{dost}, once more
  \item \emph{Say it:} \emph{śũ te tamāro dost che?} — Is he your pal?
  \item \emph{Recall:} say the Gujarati for nephew, then the Gujarati for neighbour, then say \emph{dost} again
\end{itemize}

\begin{tcolorbox}[breakable,colback=teal!4,colframe=teal!35!black,title={Before you move on}]
What does \gu{દોસ્ત} mean? (\textquotedblleft{}Pal, close friend\textquotedblright{}.) Say it once more.
\end{tcolorbox}
\section[\texorpdfstring{\gu{બહેનપણી} (bahenpaṇī)}{બહેનપણી (bahenpaṇī)}]{\gu{બહેનપણી} (bahenpaṇī) — female friend, a girl's friend}

\label{lesson:GU-C264-bahenpani}



\begin{quote}
Before the new one: say the Gujarati for neighbour, then the Gujarati for pal.
\end{quote}

\subsection*{You'll want to know: \gu{બહેનપણી}}
\textbf{\gu{બહેનપણી}} — \emph{bahenpaṇī} — \textquotedblleft{}female friend, a girl's friend\textquotedblright{}.

\begin{quote}
\textbf{\gu{મારી} \gu{બહેનપણી} \gu{આજે} \gu{નથી} \gu{આવતી}.} — \emph{mārī bahenpaṇī āje nathī āvtī} — \textquotedblleft{}My friend is not coming today.\textquotedblright{}
\end{quote}

A girl's or a woman's female friend.

\begin{cousinweb}
\textbf{\gu{બહેનપણી}} (\emph{bahenpaṇī}) is built on the taught \emph{bahen}, \textquotedblleft{}sister\textquotedblright{}, with \emph{-paṇī}. CDIAL derives \emph{bahen} from Sanskrit \emph{bhaginī}, \textquotedblleft{}sister\textquotedblright{} (CDIAL 9349).
\end{cousinweb}

\begin{grammarlens}[title={What you've built}]
One more word for everyday life at A2.
\end{grammarlens}

\subsection*{Your turn}
\begin{itemize}
\raggedright
  \item \emph{Say it:} \emph{bahenpaṇī}
  \item \emph{Say it:} \emph{bahenpaṇī}, once more
  \item \emph{Say it:} \emph{mārī bahenpaṇī āje nathī āvtī} — My friend is not coming today.
  \item \emph{Recall:} say the Gujarati for neighbour, then the Gujarati for pal, then say \emph{bahenpaṇī} again
\end{itemize}

\begin{tcolorbox}[breakable,colback=teal!4,colframe=teal!35!black,title={Before you move on}]
What does \gu{બહેનપણી} mean? (\textquotedblleft{}Female friend, a girl's friend\textquotedblright{}.) Say it once more.
\end{tcolorbox}
\section[\texorpdfstring{\gu{ઓળખવું} (oḷakhvũ)}{ઓળખવું (oḷakhvũ)}]{\gu{ઓળખવું} (oḷakhvũ) — to recognize, to know someone}

\label{lesson:GU-C264-olakhvu}



\begin{quote}
Before the new one: say the Gujarati for pal, then the Gujarati for female friend.
\end{quote}

\subsection*{You'll want to know: \gu{ઓળખવું}}
\textbf{\gu{ઓળખવું}} — \emph{oḷakhvũ} — \textquotedblleft{}to recognize, to know someone\textquotedblright{}.

\begin{quote}
\textbf{\gu{શું} \gu{તમે} \gu{મને} \gu{ઓળખો} \gu{છો}?} — \emph{śũ tame mane oḷakho cho?} — \textquotedblleft{}Do you recognize me?\textquotedblright{}
\end{quote}

\begin{cousinweb}
\textbf{\gu{ઓળખવું}} (\emph{oḷakhvũ}) is a tadbhava. It comes from Sanskrit \emph{upalakṣayati}, \textquotedblleft{}beholds, characterizes\textquotedblright{}; CDIAL glosses the Gujarati verb \textquotedblleft{}to know, distinguish, have friendship with\textquotedblright{} (CDIAL 2226).
\end{cousinweb}

\begin{grammarlens}[title={What you've built}]
One more word for everyday life at A2.
\end{grammarlens}

\subsection*{Your turn}
\begin{itemize}
\raggedright
  \item \emph{Say it:} \emph{oḷakhvũ}
  \item \emph{Say it:} \emph{oḷakhvũ}, once more
  \item \emph{Say it:} \emph{śũ tame mane oḷakho cho?} — Do you recognize me?
  \item \emph{Recall:} say the Gujarati for pal, then the Gujarati for female friend, then say \emph{oḷakhvũ} again
\end{itemize}

\begin{tcolorbox}[breakable,colback=teal!4,colframe=teal!35!black,title={Before you move on}]
What does \gu{ઓળખવું} mean? (\textquotedblleft{}To recognize, to know someone\textquotedblright{}.) Say it once more.
\end{tcolorbox}
\section[\texorpdfstring{\gu{ભેટવું} (bheṭvũ)}{ભેટવું (bheṭvũ)}]{\gu{ભેટવું} (bheṭvũ) — to embrace, to hug}

\label{lesson:GU-C264-bhetvu}



\begin{quote}
Before the new one: say the Gujarati for female friend, then the Gujarati for to recognize.
\end{quote}

\subsection*{You'll want to know: \gu{ભેટવું}}
\textbf{\gu{ભેટવું}} — \emph{bheṭvũ} — \textquotedblleft{}to embrace, to hug\textquotedblright{}.

\begin{quote}
\textbf{\gu{શું} \gu{તે} \gu{તેની} \gu{માતાને} \gu{ભેટ્યો}?} — \emph{śũ te tenī mātāne bheṭyo?} — \textquotedblleft{}Did he hug his mother?\textquotedblright{}
\end{quote}

The taught noun \textbf{\gu{ભેટ}}, \emph{bheṭ}, means a gift; the verb means to embrace.

\begin{cousinweb}
\textbf{\gu{ભેટવું}} (\emph{bheṭvũ}) is a tadbhava. It comes from \emph{*bhēṭṭ-}, a reconstructed form listed under \emph{*bhiṭ}, \textquotedblleft{}meet, throng\textquotedblright{}; CDIAL glosses the Gujarati verb \textquotedblleft{}to meet\textquotedblright{} (CDIAL 9490).
\end{cousinweb}

\begin{grammarlens}[title={What you've built}]
One more word for everyday life at A2.
\end{grammarlens}

\subsection*{Your turn}
\begin{itemize}
\raggedright
  \item \emph{Say it:} \emph{bheṭvũ}
  \item \emph{Say it:} \emph{bheṭvũ}, once more
  \item \emph{Say it:} \emph{śũ te tenī mātāne bheṭyo?} — Did he hug his mother?
  \item \emph{Recall:} say the Gujarati for female friend, then the Gujarati for to recognize, then say \emph{bheṭvũ} again
\end{itemize}

\begin{tcolorbox}[breakable,colback=teal!4,colframe=teal!35!black,title={Before you move on}]
What does \gu{ભેટવું} mean? (\textquotedblleft{}To embrace, to hug\textquotedblright{}.) Say it once more.
\end{tcolorbox}
